\begin{filecontents*}[overwrite]{\jobname.ist}
headings_flag 1
numhead_positive "符号"
symhead_positive "符号"
heading_prefix "{\\bfseries "
heading_suffix "}\\nopagebreak\\par\\vskip 2pt "
delim_0 "，"
delim_1 "，"
delim_2 "，"
delim_n "，"
\end{filecontents*}


\usepackage[
  paperwidth=170mm,
  paperheight=248.2mm,
  inner=22mm,
  outer=18mm,
  top=20mm,
  bottom=22mm,
  headheight=15pt
]{geometry}
\usepackage{amsmath,amssymb,mathtools,amsthm,bm}
\usepackage{graphicx}
\usepackage{pict2e}
\usepackage{booktabs,longtable,array}
\usepackage{enumitem}
\usepackage{multicol}
\usepackage{imakeidx}
\usepackage{xcolor}
\usepackage{fancyhdr}
\usepackage{microtype}
\usepackage{hyperref}

\definecolor{liftblue}{HTML}{1F4E79}
\hypersetup{
  unicode=true,
  colorlinks=true,
  linkcolor=liftblue,
  citecolor=liftblue,
  urlcolor=liftblue,
  pdftitle={拉格朗日交弗洛尔理论（GLM \& GPT 译）},
  pdfauthor={深谷贤治; 吴龙根; 太田啓史; 小野薫}
}

\ctexset{
  chapter = {
    name = {第,章},
    number = \arabic{chapter}
  }
}
\setcounter{secnumdepth}{3}
\setcounter{tocdepth}{2}
\setlength{\parindent}{2em}
\setlength{\parskip}{0pt}
\emergencystretch=2em
\clubpenalty=10000
\widowpenalty=10000
\displaywidowpenalty=10000
\renewcommand{\contentsname}{目录}
\renewcommand{\figurename}{图}
\renewcommand{\tablename}{表}
\renewcommand{\proofname}{证明}
\newcommand{\translatornote}[1]{\footnote{\textup{译者注：#1}}}

\makeindex[title={主题索引},noautomatic]
\expandafter\def\csname liftmarkI001\endcsname{\index{anuoerdecai xiang@阿诺尔德猜想（Arnold's conjecture）}}
\expandafter\def\csname liftmarkI002\endcsname{\index{awuqiong daishu@$A_\infty$ 代数（$A_\infty$ algebra）}}
\expandafter\def\csname liftmarkI003\endcsname{\index{awuqiong daishu@$A_\infty$ 代数（$A_\infty$ algebra）!awuqiong gongshi@$A_\infty$ 公式（$A_\infty$ formula）}}
\expandafter\def\csname liftmarkI004\endcsname{\index{awuqiong daishu@$A_\infty$ 代数（$A_\infty$ algebra）!dingjie yulian@定界余链（bounding cochain）}}
\expandafter\def\csname liftmarkI005\endcsname{\index{awuqiong daishu@$A_\infty$ 代数（$A_\infty$ algebra）!pu xulie@谱序列（spectral sequence）}}
\expandafter\def\csname liftmarkI006\endcsname{\index{awuqiong fanchou@$A_\infty$ 范畴（$A_\infty$ category）}}
\expandafter\def\csname liftmarkI007\endcsname{\index{b chang@B-场（B-field）}}
\expandafter\def\csname liftrangestartI008\endcsname{\index{bar fuxing@bar 复形（bar complex）|(}}
\expandafter\def\csname liftrangeendI008\endcsname{\index{bar fuxing@bar 复形（bar complex）|)}}
\expandafter\def\csname liftmarkI009\endcsname{\index{dingjie yulian@定界余链（bounding cochain）}}
\expandafter\def\csname liftmarkI010\endcsname{\index{bulk xingbian@bulk 形变（bulk deformations）}}
\expandafter\def\csname liftmarkI011\endcsname{\index{kalabi qiu liuxing@卡拉比--丘流形（Calabi--Yau manifolds）}}
\expandafter\def\csname liftmarkI012\endcsname{\index{kalabi qiu liuxing@卡拉比--丘流形（Calabi--Yau manifolds）!jida tuihua zu@极大退化族（maximally degenerating family）}}
\expandafter\def\csname liftmarkI013\endcsname{\index{ningjuceng daochu fanchou@凝聚层导出范畴（derived category of coherent sheaves）}}
\expandafter\def\csname liftmarkI014\endcsname{\index{fuluoer shang tongdiao@弗洛尔（上）同调（Floer (co)homology）}}
\expandafter\def\csname liftmarkI015\endcsname{\index{fuluoer shang tongdiao@弗洛尔（上）同调（Floer (co)homology）!lagelangri zi liuxing dui@拉格朗日子流形对（pair of Lagrangian submanifolds）}}
\expandafter\def\csname liftmarkI016\endcsname{\index{fuluoer shang tongdiao@弗洛尔（上）同调（Floer (co)homology）!nao bufen@挠部分（torsion part）}}
\expandafter\def\csname liftmarkI017\endcsname{\index{xingshi gaixing@形式概形（formal scheme）}}
\expandafter\def\csname liftmarkI018\endcsname{\index{xingshi gaixing@形式概形（formal scheme）!xingshi yingshe@形式映射（formal map）}}
\expandafter\def\csname liftmarkI019\endcsname{\index{tongdiao jingxiang duicheng@同调镜像对称（homological mirror symmetry）}}
\expandafter\def\csname liftmarkI020\endcsname{\index{lagelangri zi liuxing@拉格朗日子流形（Lagrangian submanifold）}}
\expandafter\def\csname liftmarkI021\endcsname{\index{lagelangri zi liuxing@拉格朗日子流形（Lagrangian submanifold）!youli lagelangri zi liuxing@有理拉格朗日子流形（rational Lagrangian submanifold）}}
\expandafter\def\csname liftmarkI022\endcsname{\index{masiluofu lei caixiang@马斯洛夫类猜想（Maslov class conjecture）}}
\expandafter\def\csname liftmarkI023\endcsname{\index{masiluofu zhibiao@马斯洛夫指标（Maslov index）!bote moersi banben@博特--莫尔斯版本（Bott--Morse version）}}
\expandafter\def\csname liftmarkI024\endcsname{\index{nuoweikefu huan@诺维科夫环（Novikov ring）!wanyou nuoweikefu huan@万有诺维科夫环（universal Novikov ring）}}
\expandafter\def\csname liftmarkI025\endcsname{\index{liangzi beiji@量子杯积（quantum cup-product）}}
\expandafter\def\csname liftmarkI026\endcsname{\index{xiangdui spin@相对 spin（relatively spin）}}
\expandafter\def\csname liftmarkI027\endcsname{\index{xiangdui spin dui@相对 spin 对（relatively spin pair）}}
\expandafter\def\csname liftmarkI028\endcsname{\index{yiwei cishu@移位次数（shifted degree）}}
\expandafter\def\csname liftmarkI029\endcsname{\index{pu xulie@谱序列（spectral sequence）!tuihua@退化（degeneration）}}
\expandafter\def\csname liftmarkI030\endcsname{\index{siteluominge qiu chengtong zhasiluo gouxiang@斯特罗明格--丘成桐--扎斯洛构想（Strominger--Yau--Zaslow proposal）}}
\expandafter\def\csname liftmarkI031\endcsname{\index{huanmian liuxing@环面流形（toric manifolds）}}
\expandafter\def\csname liftmarkI032\endcsname{\index{huanmian liuxing@环面流形（toric manifolds）!pingheng huanmian xianwei@平衡环面纤维（balanced torus fibers）}}
\expandafter\def\csname liftmarkI033\endcsname{\index{nao@挠（torsion）}}
\expandafter\def\csname liftmarkI034\endcsname{\index{tangchuan ouhe@汤川耦合（Yukawa coupling）}}
\expandafter\def\csname liftmarkI035\endcsname{\index{yuji@余积（coproduct）}}
\expandafter\def\csname liftmarkI036\endcsname{\index{delamu@德拉姆（de Rham）!fuxing@复形（complex）}}
\expandafter\def\csname liftmarkI037\endcsname{\index{cangxi jiegou@仓西结构（Kuranishi structure）}}
\expandafter\def\csname liftmarkI038\endcsname{\index{suanzi q@算子 $\mathfrak q$（operator $\mathfrak q$）}}
\expandafter\def\csname liftmarkI039\endcsname{\index{zuoyongliang fanhan@作用量泛函（action functional）}}
\expandafter\def\csname liftmarkI040\endcsname{\index{zuoyongliang fanhan@作用量泛函（action functional）!lianjie guidao@连接轨道（connecting orbits）}}
\expandafter\def\csname liftmarkI041\endcsname{\index{zuoyongliang fanhan@作用量泛函（action functional）!l2 tidu fangcheng@$L^2$-梯度方程（$L^2$-gradient equation）}}
\expandafter\def\csname liftrangestartI042\endcsname{\index{bote moersi lilun@博特--莫尔斯理论（Bott--Morse theory）|(}}
\expandafter\def\csname liftrangeendI042\endcsname{\index{bote moersi lilun@博特--莫尔斯理论（Bott--Morse theory）|)}}
\expandafter\def\csname liftmarkI043\endcsname{\index{cong dui@丛对（bundle pair）}}
\expandafter\def\csname liftmarkI044\endcsname{\index{kelifude huanmian@克利福德环面（Clifford torus）}}
\expandafter\def\csname liftrangestartI045\endcsname{\index{jishu wenti@计数问题（counting problem）|(}}
\expandafter\def\csname liftrangeendI045\endcsname{\index{jishu wenti@计数问题（counting problem）|)}}
\expandafter\def\csname liftmarkI046\endcsname{\index{fuluoer shanglian mo@弗洛尔（上）链模（Floer (co)chain module）}}
\expandafter\def\csname liftmarkI047\endcsname{\index{fuluoer shang tongdiao@弗洛尔（上）同调（Floer (co)homology）!fanchang chuxian@反常出现（anomaly appearance）}}
\expandafter\def\csname liftmarkI048\endcsname{\index{fuluoer shang tongdiao@弗洛尔（上）同调（Floer (co)homology）!bote moersi banben@博特--莫尔斯版本（Bott--Morse version）}}
\expandafter\def\csname liftmarkI049\endcsname{\index{fuluoer shang tongdiao@弗洛尔（上）同调（Floer (co)homology）!z2 qingxing xia de dingyi@$\mathbb Z_2$ 情形下的定义（definition in $\mathbb Z_2$-case）}}
\expandafter\def\csname liftmarkI050\endcsname{\index{fuluoer shang tongdiao@弗洛尔（上）同调（Floer (co)homology）!pu xulie@谱序列（spectral sequence）}}
\expandafter\def\csname liftmarkI051\endcsname{\index{fuluoer mo kongjian@弗洛尔模空间（Floer moduli space）}}
\expandafter\def\csname liftmarkI052\endcsname{\index{geluomofu jinxing dingli@格罗莫夫紧性定理（Gromov's compactness theorem）}}
\expandafter\def\csname liftmarkI053\endcsname{\index{lagelangri zi liuxing@拉格朗日子流形（Lagrangian submanifold）!dandiao lagelangri zi liuxing@单调拉格朗日子流形（monotone Lagrangian submanifolds）}}
\expandafter\def\csname liftmarkI054\endcsname{\index{masiluofu zhibiao@马斯洛夫指标（Maslov index）}}
\expandafter\def\csname liftmarkI055\endcsname{\index{masiluofu zhibiao@马斯洛夫指标（Maslov index）!lagelangri gelasimann liuxing@拉格朗日格拉斯曼流形（Lagrangian Grassmannian）}}
\expandafter\def\csname liftmarkI056\endcsname{\index{masiluofu zhibiao@马斯洛夫指标（Maslov index）!masiluofu moersi zhibiao@马斯洛夫--莫尔斯指标（Maslov--Morse index）}}
\expandafter\def\csname liftmarkI057\endcsname{\index{masiluofu zhibiao@马斯洛夫指标（Maslov index）!jixiao masiluofu shu@极小马斯洛夫数（minimal Maslov number）}}
\expandafter\def\csname liftmarkI058\endcsname{\index{masiluofu zhibiao@马斯洛夫指标（Maslov index）!quan shi zi kongjian@全实子空间（totally real subspace）}}
\expandafter\def\csname liftrangestartI059\endcsname{\index{youbian wending yingshe mo kongjian@有边稳定映射模空间（moduli space of bordered stable maps）|(}}
\expandafter\def\csname liftrangeendI059\endcsname{\index{youbian wending yingshe mo kongjian@有边稳定映射模空间（moduli space of bordered stable maps）|)}}
\expandafter\def\csname liftmarkI060\endcsname{\index{youbian wending yingshe mo kongjian@有边稳定映射模空间（moduli space of bordered stable maps）!kuisun ling wending yingshe@亏格零稳定映射（genus zero stable map）}}
\expandafter\def\csname liftmarkI061\endcsname{\index{youbian wending yingshe mo kongjian@有边稳定映射模空间（moduli space of bordered stable maps）!zhu fenzhi@主分支（main component）}}
\expandafter\def\csname liftmarkI062\endcsname{\index{youbian wending yingshe mo kongjian@有边稳定映射模空间（moduli space of bordered stable maps）!dai shi biaoji dian de banwending quxian@带实标记点的半稳定曲线（real marked semi-stable curve）}}
\expandafter\def\csname liftmarkI063\endcsname{\index{nuoweikefu fugai kongjian@诺维科夫覆盖空间（Novikov covering space）}}
\expandafter\def\csname liftmarkI064\endcsname{\index{nuoweikefu huan@诺维科夫环（Novikov ring）}}
\expandafter\def\csname liftmarkI065\endcsname{\index{zhangai@障碍（obstruction）!shouyao zhangai@首要障碍（primary obstruction）}}
\expandafter\def\csname liftmarkI066\endcsname{\index{zai mou chu danshe@在某处单射（somewhere injective）}}
\expandafter\def\csname liftmarkI067\endcsname{\index{awuqiong daishu@$A_\infty$ 代数（$A_\infty$ algebra）!gyouxi@$G$-有隙（$G$-gapped）}}
\expandafter\def\csname liftmarkI068\endcsname{\index{awuqiong daishu@$A_\infty$ 代数（$A_\infty$ algebra）!maurer cartan fangcheng@Maurer--Cartan 方程（Maurer--Cartan equation）}}
\expandafter\def\csname liftmarkI069\endcsname{\index{awuqiong daishu@$A_\infty$ 代数（$A_\infty$ algebra）!wuzhangai@无障碍（unobstructed）}}
\expandafter\def\csname liftmarkI070\endcsname{\index{awuqiong daishu@$A_\infty$ 代数（$A_\infty$ algebra）!ruowuzhangai@弱无障碍（weakly unobstructed）}}
\expandafter\def\csname liftmarkI071\endcsname{\index{awuqiong daishu@$A_\infty$ 代数（$A_\infty$ algebra）!bulk xingbian hou ruowuzhangai@经 bulk 形变后（弱）无障碍（(weakly) unobstructed after bulk deformations）}}
\expandafter\def\csname liftmarkI072\endcsname{\index{bulk xingbian@bulk 形变（bulk deformations）!bulk xingbian hou ruowuzhangai@经 bulk 形变后（弱）无障碍（(weakly) unobstructed after bulk deformation）}}
\expandafter\def\csname liftrangestartI073\endcsname{\index{fuluoer shang tongdiao@弗洛尔（上）同调（Floer (co)homology）!zhangai lei@障碍类（obstruction classes）|(}}
\expandafter\def\csname liftrangeendI073\endcsname{\index{fuluoer shang tongdiao@弗洛尔（上）同调（Floer (co)homology）!zhangai lei@障碍类（obstruction classes）|)}}
\expandafter\def\csname liftrangestartI074\endcsname{\index{geluomofu weiteng bubianliang@格罗莫夫--威腾不变量（Gromov--Witten invariant）|(}}
\expandafter\def\csname liftrangeendI074\endcsname{\index{geluomofu weiteng bubianliang@格罗莫夫--威腾不变量（Gromov--Witten invariant）|)}}
\expandafter\def\csname liftrangestartI075\endcsname{\index{lagelangri zi liuxing@拉格朗日子流形（Lagrangian submanifold）!wuzhangai@无障碍（unobstructed）|(}}
\expandafter\def\csname liftrangeendI075\endcsname{\index{lagelangri zi liuxing@拉格朗日子流形（Lagrangian submanifold）!wuzhangai@无障碍（unobstructed）|)}}
\expandafter\def\csname liftmarkI076\endcsname{\index{lagelangri zi liuxing@拉格朗日子流形（Lagrangian submanifold）!ruowuzhangai@弱无障碍（weakly unobstructed）}}
\expandafter\def\csname liftmarkI077\endcsname{\index{lagelangri zi liuxing@拉格朗日子流形（Lagrangian submanifold）!bulk xingbian hou ruowuzhangai@经 bulk 形变后（弱）无障碍（(weakly) unobstructed after bulk deformations）}}
\expandafter\def\csname liftmarkI078\endcsname{\index{maurer cartan fangcheng@Maurer--Cartan 方程（Maurer--Cartan equation）}}
\expandafter\def\csname liftmarkI079\endcsname{\index{shi hanshu@势函数（potential function）}}
\expandafter\def\csname liftmarkI080\endcsname{\index{awuqiong daishu@$A_\infty$ 代数（$A_\infty$ algebra）!xingbian@形变（deformation）}}
\expandafter\def\csname liftmarkI081\endcsname{\index{awuqiong daishu@$A_\infty$ 代数（$A_\infty$ algebra）!nengliang lvguo@能量滤过（energy filtration）}}
\expandafter\def\csname liftmarkI082\endcsname{\index{awuqiong daishu@$A_\infty$ 代数（$A_\infty$ algebra）!lvguo awuqiong daishu@滤过 $A_\infty$ 代数（filtered $A_\infty$ algebra）}}
\expandafter\def\csname liftmarkI083\endcsname{\index{awuqiong daishu@$A_\infty$ 代数（$A_\infty$ algebra）!gyouxi lvguo awuqiong daishu@$G$-有隙滤过 $A_\infty$ 代数（$G$-gapped filtered $A_\infty$ algebras）}}
\expandafter\def\csname liftmarkI084\endcsname{\index{awuqiong daishu@$A_\infty$ 代数（$A_\infty$ algebra）!r yuehua@$R$-约化（$R$-reduction）}}
\expandafter\def\csname liftmarkI085\endcsname{\index{awuqiong daishu@$A_\infty$ 代数（$A_\infty$ algebra）!yange@严格（strict）}}
\expandafter\def\csname liftmarkI086\endcsname{\index{awuqiong daishu@$A_\infty$ 代数（$A_\infty$ algebra）!yaoyuan@幺元（unit）}}
\expandafter\def\csname liftmarkI087\endcsname{\index{awuqiong daishu tongtai@$A_\infty$（代数）同态（$A_\infty$ (algebra) homomorphism）}}
\expandafter\def\csname liftmarkI088\endcsname{\index{awuqiong daishu tongtai@$A_\infty$（代数）同态（$A_\infty$ (algebra) homomorphism）!fuhe@复合（composition）}}
\expandafter\def\csname liftmarkI089\endcsname{\index{awuqiong daishu tongtai@$A_\infty$（代数）同态（$A_\infty$ (algebra) homomorphism）!lvguo awuqiong tongtai@滤过 $A_\infty$ 同态（filtered $A_\infty$ homomorphisms）}}
\expandafter\def\csname liftmarkI090\endcsname{\index{awuqiong daishu tongtai@$A_\infty$（代数）同态（$A_\infty$ (algebra) homomorphism）!gyouxi@$G$-有隙（$G$-gapped）}}
\expandafter\def\csname liftmarkI091\endcsname{\index{awuqiong daishu tongtai@$A_\infty$（代数）同态（$A_\infty$ (algebra) homomorphism）!yange@严格（strict）}}
\expandafter\def\csname liftmarkI092\endcsname{\index{yudaozi@余导子（coderivation）}}
\expandafter\def\csname liftmarkI093\endcsname{\index{nengliang lvguo@能量滤过（energy filtration）}}
\expandafter\def\csname liftmarkI094\endcsname{\index{shumu lvguo@数目滤过（number filtration）}}
\expandafter\def\csname liftmarkI095\endcsname{\index{bar fuxing@bar 复形（bar complex）}}
\expandafter\def\csname liftmarkI096\endcsname{\index{bar fuxing@bar 复形（bar complex）!duichenghua@对称化（symmetrization）}}
\expandafter\def\csname liftrangestartI097\endcsname{\index{lwuqiong daishu@$L_\infty$ 代数（$L_\infty$ algebra）|(}}
\expandafter\def\csname liftrangeendI097\endcsname{\index{lwuqiong daishu@$L_\infty$ 代数（$L_\infty$ algebra）|)}}
\expandafter\def\csname liftmarkI098\endcsname{\index{awuqiong daishu tonglun dengjia@$A_\infty$ 代数的同伦等价（homotopy equivalence of $A_\infty$ algebra）!ruo tonglun dengjia@弱同伦等价（weak homotopy equivalence）}}
\expandafter\def\csname liftrangestartI099\endcsname{\index{awuqiong daishu@$A_\infty$ 代数（$A_\infty$ algebra）!tonglun yaoyuan@同伦幺元（homotopy unit）|(}}
\expandafter\def\csname liftrangeendI099\endcsname{\index{awuqiong daishu@$A_\infty$ 代数（$A_\infty$ algebra）!tonglun yaoyuan@同伦幺元（homotopy unit）|)}}
\expandafter\def\csname liftmarkI100\endcsname{\index{beiji de awuqiong xingbian@杯积的 $A_\infty$ 形变（$A_\infty$ deformation of cup product）}}
\expandafter\def\csname liftmarkI101\endcsname{\index{awuqiong daishu tongtai@$A_\infty$（代数）同态（$A_\infty$ (algebra) homomorphism）!hanyao tonglun hanyao@含幺（相应地，同伦含幺）（unital (resp. homotopy unital)）}}
\expandafter\def\csname liftmarkI102\endcsname{\index{tonglun yaoyuan@同伦幺元（homotopy unit）}}
\expandafter\def\csname liftmarkI103\endcsname{\index{daizhuang shu@带状树（ribbon tree）}}
\expandafter\def\csname liftrangestartI104\endcsname{\index{awuqiong daishu@$A_\infty$ 代数（$A_\infty$ algebra）!lvguo awuqiong daishu guanlian lagelangri zi liuxing@与拉格朗日子流形关联的滤过 $A_\infty$ 代数（filtered $A_\infty$ algebra associated to Lagrangian submanifolds）|(}}
\expandafter\def\csname liftrangeendI104\endcsname{\index{awuqiong daishu@$A_\infty$ 代数（$A_\infty$ algebra）!lvguo awuqiong daishu guanlian lagelangri zi liuxing@与拉格朗日子流形关联的滤过 $A_\infty$ 代数（filtered $A_\infty$ algebra associated to Lagrangian submanifolds）|)}}
\expandafter\def\csname liftmarkI105\endcsname{\index{xu jiben lian@虚基本链（virtual fundamental chain）}}
\expandafter\def\csname liftmarkI106\endcsname{\index{awuqiong daishu@$A_\infty$ 代数（$A_\infty$ algebra）!fuluoer shang tongdiao@弗洛尔上同调（Floer cohomology）}}
\expandafter\def\csname liftmarkI107\endcsname{\index{awuqiong daishu@$A_\infty$ 代数（$A_\infty$ algebra）!zhangai@障碍（obstructions）}}
\expandafter\def\csname liftmarkI108\endcsname{\index{fuluoer shang tongdiao@弗洛尔（上）同调（Floer (co)homology）!zhangai lei@障碍类（obstruction classes）}}
\expandafter\def\csname liftmarkI109\endcsname{\index{lagelangri zi liuxing@拉格朗日子流形（Lagrangian submanifold）!wuzhangai@无障碍（unobstructed）}}
\expandafter\def\csname liftmarkI110\endcsname{\index{zhangai@障碍（obstruction）!wuzhangai awuqiong daishu@无障碍 $A_\infty$ 代数（unobstructed $A_\infty$ algebras）}}
\expandafter\def\csname liftmarkI111\endcsname{\index{dingjie yulian@定界余链（bounding cochain）!ruo dingjie yulian@弱定界余链（weak bounding cochain）}}
\expandafter\def\csname liftmarkI112\endcsname{\index{awuqiong daishu@$A_\infty$ 代数（$A_\infty$ algebra）!zai g shang ruowuzhangai@在 $G$ 上弱无障碍（weakly unobstructed over $G$）}}
\expandafter\def\csname liftmarkI113\endcsname{\index{xunhuan@循环（Cyclic）!xunhuan duichengxing@循环对称性（symmetry）}}
\expandafter\def\csname liftrangestartI114\endcsname{\index{chaoshi@超势（superpotential）|(}}
\expandafter\def\csname liftrangeendI114\endcsname{\index{chaoshi@超势（superpotential）|)}}
\expandafter\def\csname liftmarkI115\endcsname{\index{awuqiong shuangmo@$A_\infty$ 双模（$A_\infty$ bimodule）}}
\expandafter\def\csname liftmarkI116\endcsname{\index{awuqiong shuangmo@$A_\infty$ 双模（$A_\infty$ bimodule）!lvguo awuqiong shuangmo@滤过 $A_\infty$ 双模（filtered $A_\infty$ bimodule）}}
\expandafter\def\csname liftmarkI117\endcsname{\index{awuqiong shuangmo@$A_\infty$ 双模（$A_\infty$ bimodule）!r yuehua@$R$-约化（$R$-reduction）}}
\expandafter\def\csname liftmarkI118\endcsname{\index{awuqiong shuangmo@$A_\infty$ 双模（$A_\infty$ bimodule）!hanyao tonglun hanyao@含幺、同伦含幺（unital, homotopy unital）}}
\expandafter\def\csname liftmarkI119\endcsname{\index{awuqiong shuangmo tongtai@$A_\infty$ 双模同态（$A_\infty$ bimodule homomorphism）}}
\expandafter\def\csname liftmarkI120\endcsname{\index{awuqiong shuangmo tongtai@$A_\infty$ 双模同态（$A_\infty$ bimodule homomorphism）!gyouxi@$G$-有隙（$G$-gapped）}}
\expandafter\def\csname liftrangestartI121\endcsname{\index{awuqiong shuangmo tongtai@$A_\infty$ 双模同态（$A_\infty$ bimodule homomorphism）!hanyao tonglun hanyao lvguo awuqiong shuangmo tongtai@含幺、同伦含幺滤过 $A_\infty$ 双模同态（unital, homotopy unital filtered $A_\infty$ bimodule homomorphism）|(}}
\expandafter\def\csname liftrangeendI121\endcsname{\index{awuqiong shuangmo tongtai@$A_\infty$ 双模同态（$A_\infty$ bimodule homomorphism）!hanyao tonglun hanyao lvguo awuqiong shuangmo tongtai@含幺、同伦含幺滤过 $A_\infty$ 双模同态（unital, homotopy unital filtered $A_\infty$ bimodule homomorphism）|)}}
\expandafter\def\csname liftmarkI122\endcsname{\index{awuqiong shuangmo@$A_\infty$ 双模（$A_\infty$ bimodule）!xingbian awuqiong shuangmo@形变 $A_\infty$ 双模（deformed $A_\infty$ bimodule）}}
\expandafter\def\csname liftmarkI123\endcsname{\index{awuqiong shuangmo@$A_\infty$ 双模（$A_\infty$ bimodule）!gyouxi lvguo@$G$-有隙滤过（$G$-gapped filtered）}}
\expandafter\def\csname liftmarkI124\endcsname{\index{yudaozi@余导子（coderivation）!shuang yudaozi@双余导子（bi-coderivation）}}
\expandafter\def\csname liftmarkI125\endcsname{\index{awuqiong shuangmo@$A_\infty$ 双模（$A_\infty$ bimodule）!fuluoer shang tongdiao@弗洛尔上同调（Floer cohomology）}}
\expandafter\def\csname liftrangestartI126\endcsname{\index{fuluoer mo kongjian@弗洛尔模空间（Floer moduli space）!jinhua@紧化（compactification）|(}}
\expandafter\def\csname liftrangeendI126\endcsname{\index{fuluoer mo kongjian@弗洛尔模空间（Floer moduli space）!jinhua@紧化（compactification）|)}}
\expandafter\def\csname liftmarkI127\endcsname{\index{fuluoer mo kongjian@弗洛尔模空间（Floer moduli space）!dai biaoji wending duanlie guiji@带标记稳定断裂轨迹（marked stable broken trajectories）}}
\expandafter\def\csname liftmarkI128\endcsname{\index{fuluoer mo kongjian@弗洛尔模空间（Floer moduli space）!zi yaobanqun g l1 l0@子幺半群 $G(L^{(1)},L^{(0)})$（submonoid $G(L^{(1)},L^{(0)})$）}}
\expandafter\def\csname liftrangestartI129\endcsname{\index{masiluofu zhibiao@马斯洛夫指标（Maslov index）!bote moersi banben@博特--莫尔斯版本（Bott--Morse version）|(}}
\expandafter\def\csname liftrangeendI129\endcsname{\index{masiluofu zhibiao@马斯洛夫指标（Maslov index）!bote moersi banben@博特--莫尔斯版本（Bott--Morse version）|)}}
\expandafter\def\csname liftmarkI130\endcsname{\index{fuluoer shang tongdiao@弗洛尔（上）同调（Floer (co)homology）!chengfa jiegou@乘法结构（multiplicative structures）}}
\expandafter\def\csname liftmarkI131\endcsname{\index{masiluofu zhibiao@马斯洛夫指标（Maslov index）!bote moersi banben@博特--莫尔斯版本（Bott--Morse version）}}
\expandafter\def\csname liftmarkI132\endcsname{\index{bar fuxing@bar 复形（bar complex）!xunhuan bar fuxing@循环 bar 复形（cyclic bar complex）}}
\expandafter\def\csname liftmarkI133\endcsname{\index{xunhuan@循环（Cyclic）!bar fuxing@bar 复形（bar complex）}}
\expandafter\def\csname liftmarkI134\endcsname{\index{geluomofu weiteng bubianliang@格罗莫夫--威腾不变量（Gromov--Witten invariant）}}
\expandafter\def\csname liftmarkI135\endcsname{\index{gysin tongtai@Gysin 同态（Gysin homomorphism）}}
\expandafter\def\csname liftmarkI136\endcsname{\index{suanzi p@算子 $\mathfrak p$（operator $\mathfrak p$）}}
\expandafter\def\csname liftmarkI137\endcsname{\index{suanzi p@算子 $\mathfrak p$（operator $\mathfrak p$）!xunhuan bar fuxing@循环 bar 复形（cyclic bar complex）}}
\expandafter\def\csname liftmarkI138\endcsname{\index{suanzi p@算子 $\mathfrak p$（operator $\mathfrak p$）!gw ling yi m l@$GW_{0,1}(M)(L)$}}
\expandafter\def\csname liftmarkI139\endcsname{\index{e xing@$E_*$}}
\expandafter\def\csname liftrangestartI140\endcsname{\index{fuluoer shang tongdiao@弗洛尔（上）同调（Floer (co)homology）!bubianxing@不变性（invariance）|(}}
\expandafter\def\csname liftrangeendI140\endcsname{\index{fuluoer shang tongdiao@弗洛尔（上）同调（Floer (co)homology）!bubianxing@不变性（invariance）|)}}
\expandafter\def\csname liftmarkI141\endcsname{\index{awuqiong daishu tonglun dengjia@$A_\infty$ 代数的同伦等价（homotopy equivalence of $A_\infty$ algebra）!moxing lingyi c@$[0,1]\times C$ 的模型（model of $[0,1]\times C$）}}
\expandafter\def\csname liftmarkI142\endcsname{\index{incl@$\mathrm{Incl}$}}
\expandafter\def\csname liftmarkI143\endcsname{\index{lingyi c moxing@$[0,1]\times C$ 的模型（model of $[0,1]\times C$）!awuqiong daishu@$A_\infty$ 代数（$A_\infty$ algebra）}}
\expandafter\def\csname liftmarkI144\endcsname{\index{awuqiong daishu tonglun dengjia@$A_\infty$ 代数的同伦等价（homotopy equivalence of $A_\infty$ algebra）!c lingyi mathfrak m k@$(C^{[0,1]},\mathfrak M_k)$}}
\expandafter\def\csname liftmarkI145\endcsname{\index{tonglun@同伦（homotopy）}}
\expandafter\def\csname liftmarkI146\endcsname{\index{tonglun@同伦（homotopy）!awuqiong tongtai zhijian@$A_\infty$ 同态之间（between $A_\infty$ homomorphisms）}}
\expandafter\def\csname liftmarkI147\endcsname{\index{awuqiong daishu tonglun dengjia@$A_\infty$ 代数的同伦等价（homotopy equivalence of $A_\infty$ algebra）}}
\expandafter\def\csname liftmarkI148\endcsname{\index{awuqiong daishu tonglun dengjia@$A_\infty$ 代数的同伦等价（homotopy equivalence of $A_\infty$ algebra）!tonglun hanyao de tonglun@同伦含幺地同伦（homotopy-unitally homotopic）}}
\expandafter\def\csname liftmarkI149\endcsname{\index{awuqiong daishu tonglun dengjia@$A_\infty$ 代数的同伦等价（homotopy equivalence of $A_\infty$ algebra）!yange tonglun dengjia@严格同伦等价（strict homotopy equivalence）}}
\expandafter\def\csname liftmarkI150\endcsname{\index{awuqiong daishu tonglun dengjia@$A_\infty$ 代数的同伦等价（homotopy equivalence of $A_\infty$ algebra）!hanyao tonglun@含幺同伦（unitally homotopic）}}
\expandafter\def\csname liftmarkI151\endcsname{\index{whitehead dingli@Whitehead 定理（Whitehead theorem）!awuqiong daishu@对 $A_\infty$ 代数（for $A_\infty$ algebras）}}
\expandafter\def\csname liftmarkI152\endcsname{\index{guifan dengjia@规范等价（gauge equivalence）}}
\expandafter\def\csname liftmarkI153\endcsname{\index{guifan dengjia@规范等价（gauge equivalence）!dingjie yulian@定界余链的规范等价（of bounding cochains）}}
\expandafter\def\csname liftmarkI154\endcsname{\index{guifan dengjia@规范等价（gauge equivalence）!fuluoer shang tongdiao bubianxing@弗洛尔上同调的不变性（invariance of Floer cohomology）}}
\expandafter\def\csname liftmarkI155\endcsname{\index{guifan dengjia@规范等价（gauge equivalence）!ruo dingjie yulian@弱定界余链的规范等价（of weak bounding cochains）}}
\expandafter\def\csname liftmarkI156\endcsname{\index{guifan dengjia@规范等价（gauge equivalence）!shi hanshu bubianxing@势函数的不变性（invariance of potential function）}}
\expandafter\def\csname liftmarkI157\endcsname{\index{guifan dengjia@规范等价（gauge equivalence）!b b0 b b1 zhijian@对子 $(\mathfrak b_0,b_0)$、$(\mathfrak b_1,b_1)$ 之间（between the pairs $(\mathfrak b_0,b_0)$, $(\mathfrak b_1,b_1)$）}}
\expandafter\def\csname liftmarkI158\endcsname{\index{ak daishu@$A_K$ 代数（$A_K$ algebra）}}
\expandafter\def\csname liftmarkI159\endcsname{\index{ak daishu@$A_K$ 代数（$A_K$ algebra）!ak tongtai@$A_K$ 同态（$A_K$ homomorphism）}}
\expandafter\def\csname liftmarkI160\endcsname{\index{ak daishu@$A_K$ 代数（$A_K$ algebra）!ak tonglun@$A_K$ 同伦（$A_K$ homotopy）}}
\expandafter\def\csname liftmarkI161\endcsname{\index{ak daishu@$A_K$ 代数（$A_K$ algebra）!ak tonglun dengjia@$A_K$ 同伦等价（$A_K$ homotopy equivalence）}}
\expandafter\def\csname liftmarkI162\endcsname{\index{bar fuxing@bar 复形（bar complex）!bufen yuehua bar fuxing@部分约化 bar 复形（partially reduced bar complex）}}
\expandafter\def\csname liftmarkI163\endcsname{\index{yudaozi@余导子（coderivation）}}
\expandafter\def\csname liftmarkI164\endcsname{\index{hochschild shangtongdiao@Hochschild 上同调（Hochschild cohomology）}}
\expandafter\def\csname liftmarkI165\endcsname{\index{hochschild shangtongdiao@Hochschild 上同调（Hochschild cohomology）!awuqiong tongtai@$A_\infty$ 同态（of $A_\infty$ homomorphism）}}
\expandafter\def\csname liftmarkI166\endcsname{\index{hochschild shangtongdiao@Hochschild 上同调（Hochschild cohomology）!tonglun dengjia@同伦等价（of homotopy equivalence）}}
\expandafter\def\csname liftmarkI167\endcsname{\index{nengliang lvguo@能量滤过（energy filtration）}}
\expandafter\def\csname liftmarkI168\endcsname{\index{bar fuxing@bar 复形（bar complex）!guolu bufen yuehua bar fuxing@滤过版本的部分约化 bar 复形（partially reduced bar complex for filtered version）}}
\expandafter\def\csname liftmarkI169\endcsname{\index{whitehead dingli@Whitehead 定理（Whitehead theorem）!awuqiong daishu@对 $A_\infty$ 代数（for $A_\infty$ algebras）}}
\expandafter\def\csname liftmarkI170\endcsname{\index{hochschild shangtongdiao@Hochschild 上同调（Hochschild cohomology）!awuqiong tongtai@$A_\infty$ 同态（of $A_\infty$ homomorphism）}}
\expandafter\def\csname liftmarkI171\endcsname{\index{nengliang lvguo@能量滤过（energy filtration）}}
\expandafter\def\csname liftmarkI172\endcsname{\index{awuqiong daishu tonglun dengjia@$A_\infty$ 代数的同伦等价（homotopy equivalence of $A_\infty$ algebra）!jihe shixian@几何实现（geometric realization）}}
\expandafter\def\csname liftmarkI173\endcsname{\index{shijian youxu ji@时间有序积（time ordered product, top）}}
\expandafter\def\csname liftmarkI174\endcsname{\index{zhushi ji@逐时积（time wise product, twp）}}
\expandafter\def\csname liftmarkI175\endcsname{\index{tonglun@同伦（homotopy）!awuqiong tongtai zhijian@$A_\infty$ 同态之间（between $A_\infty$ homomorphisms）}}
\expandafter\def\csname liftmarkI176\endcsname{\index{shipei qiyi danxing@适配奇异单形（adapted singular simplex）}}
\expandafter\def\csname liftmarkI177\endcsname{\index{zhushi ji@逐时积（time wise product, twp）}}
\expandafter\def\csname liftmarkI178\endcsname{\index{suanzi q@算子 $\mathfrak q$（operator $\mathfrak q$）!gaibian neibu bishanglian@改变内部闭上链（changing the cycle in the interior）}}
\expandafter\def\csname liftmarkI179\endcsname{\index{guifan dengjia@规范等价（gauge equivalence）!b b0 b b1 zhijian@对子 $(\mathfrak b_0,b_0)$、$(\mathfrak b_1,b_1)$ 之间（between the pairs $(\mathfrak b_0,b_0)$, $(\mathfrak b_1,b_1)$）}}
\expandafter\def\csname liftmarkI180\endcsname{\index{suanzi q@算子 $\mathfrak q$（operator $\mathfrak q$）!xin weifen tongpei bubianxing@辛微分同胚下的不变性（invariance under symplectic diffeomorphisms）}}
\expandafter\def\csname liftmarkI181\endcsname{\index{fuluoer shang tongdiao@弗洛尔（上）同调（Floer (co)homology）!nuoweikefu huan shang@在 $\Lambda(L^{(0)},L^{(1)};\ell_0)$ 上（over $\Lambda(L^{(0)},L^{(1)};\ell_0)$）}}
\expandafter\def\csname liftmarkI182\endcsname{\index{g jihe@$G$-集合（$G$-set）}}
\expandafter\def\csname liftmarkI183\endcsname{\index{awuqiong shuangmo tonglun dengjia@$A_\infty$ 双模的同伦等价（homotopy equivalence of $A_\infty$ bimodule）!nengliang sunshi@能量损失（energy loss）}}
\expandafter\def\csname liftmarkI184\endcsname{\index{yudaishu tonggou@余代数同构（coalgebra isomorphisms）}}
\expandafter\def\csname liftmarkI185\endcsname{\index{tonglun@同伦（homotopy）!awuqiong shuangmo tongtai zhijian@$A_\infty$ 双模同态之间（between $A_\infty$ bimodule homomorphisms）}}
\expandafter\def\csname liftmarkI186\endcsname{\index{awuqiong shuangmo tonglun dengjia@$A_\infty$ 双模的同伦等价（homotopy equivalence of $A_\infty$ bimodule）}}
\expandafter\def\csname liftmarkI187\endcsname{\index{awuqiong shuangmo tonglun dengjia@$A_\infty$ 双模的同伦等价（homotopy equivalence of $A_\infty$ bimodule）!yange tonglun dengjia@严格同伦等价（strict homotopy equivalence）}}
\expandafter\def\csname liftmarkI188\endcsname{\index{guifan dengjia@规范等价（gauge equivalence）!fuluoer shangtongdiao bubianxing@弗洛尔上同调的不变性（invariance of Floer cohomology）}}
\expandafter\def\csname liftmarkI189\endcsname{\index{akl shuangmo@$A_{K,L}$ 双模（$A_{K,L}$ bimodule）}}
\expandafter\def\csname liftmarkI190\endcsname{\index{akl shuangmo@$A_{K,L}$ 双模（$A_{K,L}$ bimodule）!tongtai@同态（homomorphism）}}
\expandafter\def\csname liftmarkI191\endcsname{\index{whitehead dingli@Whitehead 定理（Whitehead theorem）!awuqiong shuangmo@对 $A_\infty$ 双模（for $A_\infty$ bimodules）}}
\expandafter\def\csname liftmarkI192\endcsname{\index{whitehead dingli@Whitehead 定理（Whitehead theorem）!awuqiong shuangmo@对 $A_\infty$ 双模（for $A_\infty$ bimodules）}}
\expandafter\def\csname liftmarkI193\endcsname{\index{awuqiong shuangmo tonglun dengjia@$A_\infty$ 双模的同伦等价（homotopy equivalence of $A_\infty$ bimodule）}}
\expandafter\def\csname liftmarkI194\endcsname{\index{awuqiong shuangmo tonglun dengjia@$A_\infty$ 双模的同伦等价（homotopy equivalence of $A_\infty$ bimodule）!jihe shixian@几何实现（geometric realization）}}
\expandafter\def\csname liftmarkI195\endcsname{\index{tonglun@同伦（homotopy）!awuqiong shuangmo tongtai zhijian@$A_\infty$ 双模同态之间（between $A_\infty$ bimodule homomorphisms）}}
\expandafter\def\csname liftmarkI196\endcsname{\index{nengliang guji@能量估计（energy estimate）}}
\expandafter\def\csname liftmarkI197\endcsname{\index{huofeier changdu@霍费尔长度（Hofer length）}}
\expandafter\def\csname liftmarkI198\endcsname{\index{huofeier changdu@霍费尔长度（Hofer length）!hamidun tonghen@哈密顿同痕的（of Hamiltonian isotopy）}}
\expandafter\def\csname liftmarkI199\endcsname{\index{suanzi r@算子 $\mathfrak r$（operator $\mathfrak r$）}}
\expandafter\def\csname liftmarkI200\endcsname{\index{suanzi r@算子 $\mathfrak r$（operator $\mathfrak r$）!dingyi@定义（definition）}}
\expandafter\def\csname liftmarkI201\endcsname{\index{suanzi r@算子 $\mathfrak r$（operator $\mathfrak r$）!gouzao@构造（construction）}}
\expandafter\def\csname liftmarkI202\endcsname{\index{suanzi q@算子 $\mathfrak q$（operator $\mathfrak q$）!shuangmo tonglun dengjia@双模的同伦等价（homotopy equivalence of bimodule）}}
\expandafter\def\csname liftmarkI203\endcsname{\index{suanzi r@算子 $\mathfrak r$（operator $\mathfrak r$）!shuangmo tonglun dengjia@双模的同伦等价（homotopy equivalence of bimodules）}}
\expandafter\def\csname liftmarkI204\endcsname{\index{fuluoer shang tongdiao@弗洛尔（上）同调（Floer (co)homology）!bubianxing@不变性（invariance）}}
\expandafter\def\csname liftmarkI205\endcsname{\index{awuqiong shuangmo tongtai@$A_\infty$ 双模同态（$A_\infty$ bimodule homomorphism）!yange@严格（strict）}}
\expandafter\def\csname liftmarkI206\endcsname{\index{awuqiong daishu@$A_\infty$ 代数（$A_\infty$ algebra）!tonggou@同构（isomorphic）}}
\expandafter\def\csname liftmarkI207\endcsname{\index{awuqiong shuangmo@$A_\infty$ 双模（$A_\infty$ bimodule）!tonggou@同构（isomorphic）}}
\expandafter\def\csname liftmarkI208\endcsname{\index{dianfan moxing@典范模型（canonical model）}}
\expandafter\def\csname liftmarkI209\endcsname{\index{dianfan moxing@典范模型（canonical model）!lvguo awuqiong daishu@滤过 $A_\infty$ 代数的（of filtered $A_\infty$ algebra）}}
\expandafter\def\csname liftmarkI210\endcsname{\index{cangxi yingshe@仓西映射（Kuranishi map）}}
\expandafter\def\csname liftmarkI211\endcsname{\index{dianfan moxing@典范模型（canonical model）!lvguo awuqiong shuangmo@滤过 $A_\infty$ 双模的（of filtered $A_\infty$ bimodule）}}
\expandafter\def\csname liftmarkI212\endcsname{\index{ningjuceng daochu fanchou@凝聚层导出范畴（derived category of coherent sheaves）}}
\expandafter\def\csname liftmarkI213\endcsname{\index{xingshi gaixing@形式概形（formal scheme）}}
\expandafter\def\csname liftmarkI214\endcsname{\index{kedou duandian@蝌蚪端点（tad pole）}}
\expandafter\def\csname liftmarkI215\endcsname{\index{suanzi q@算子 $\mathfrak q$（operator $\mathfrak q$）}}
\expandafter\def\csname liftmarkI216\endcsname{\index{aodan wenti@奥丹问题（Audin's question）}}
\expandafter\def\csname liftrangestartI217\endcsname{\index{nao@挠（torsion）|(}}
\expandafter\def\csname liftrangeendI217\endcsname{\index{nao@挠（torsion）|)}}
\expandafter\def\csname liftmarkI218\endcsname{\index{huofeier juli@霍费尔距离（Hofer distance）}}
\expandafter\def\csname liftmarkI219\endcsname{\index{lagelangri zi liuxing@拉格朗日子流形（Lagrangian submanifold）!keweiyi lagelangri zi liuxing@可位移的拉格朗日子流形（displaceable Lagrangian submanifolds）}}
\expandafter\def\csname liftmarkI220\endcsname{\index{lagelangri zi liuxing@拉格朗日子流形（Lagrangian submanifold）!youli lagelangri zi liuxing@有理拉格朗日子流形（rational Lagrangian submanifold）}}
\expandafter\def\csname liftmarkI221\endcsname{\index{youli wuzhangai@有理无障碍（rationally unobstructed）}}
\expandafter\def\csname liftmarkI222\endcsname{\index{t rank@$t$-秩（$t$-rank）}}
\expandafter\def\csname liftmarkI223\endcsname{\index{c f z@零能生成完备自由滤过模（c.f.z.）}}
\expandafter\def\csname liftmarkI224\endcsname{\index{youxi tiaojian@有隙条件（gapped condition）}}
\expandafter\def\csname liftmarkI225\endcsname{\index{youxi shanglian yingshe@有隙上链映射（gapped cochain map）}}
\expandafter\def\csname liftmarkI226\endcsname{\index{nao@挠（torsion）}}
\expandafter\def\csname liftmarkI227\endcsname{\index{cangxi jiegou@仓西结构（Kuranishi structure）!mokongjian jiegou@$\mathcal M_{k+1}^{\mathrm{main}}(\beta)$ 上的结构（of $\mathcal M_{k+1}^{\mathrm{main}}(\beta)$）}}
\expandafter\def\csname liftmarkI228\endcsname{\index{fuleidehuomu lilun@弗雷德霍姆理论（Fredholm theory）}}
\expandafter\def\csname liftmarkI229\endcsname{\index{jiaohe@胶合（gluing）}}
\expandafter\def\csname liftmarkI230\endcsname{\index{yiwang yingshe@遗忘映射（forgetful map）}}
\expandafter\def\csname liftmarkI231\endcsname{\index{geluomofu jinxing dingli@格罗莫夫紧性定理（Gromov's compactness theorem）}}
\expandafter\def\csname liftmarkI232\endcsname{\index{geluomofu jinxing dingli@格罗莫夫紧性定理（Gromov's compactness theorem）}}
\expandafter\def\csname liftmarkI233\endcsname{\index{sitaxiefu baoqiang@斯塔谢夫胞腔（Stasheff cells）}}
\expandafter\def\csname liftmarkI234\endcsname{\index{shijian fenpei@时间分配（time allocation）}}
\expandafter\def\csname liftmarkI235\endcsname{\index{awuqiong kongjian@$A_\infty$ 空间（$A_\infty$ space）}}
\expandafter\def\csname liftmarkI236\endcsname{\index{awuqiong yingshe@$A_\infty$ 映射（$A_\infty$ map）}}
\expandafter\def\csname liftmarkI237\endcsname{\index{duijiaoxian chude hengjiexing@对角线处的横截性（transversality at diagonal）}}
\expandafter\def\csname liftmarkI238\endcsname{\index{guanghua qiyi danxingde keshuji@光滑奇异单形的可数集（countable set of smooth singular simplices）}}
\expandafter\def\csname liftmarkI239\endcsname{\index{guanghua qiyi danxingde keshuji@光滑奇异单形的可数集（countable set of smooth singular simplices）}}
\expandafter\def\csname liftmarkI240\endcsname{\index{guanghua qiyi danxingde keshuji@光滑奇异单形的可数集（countable set of smooth singular simplices）}}
\expandafter\def\csname liftmarkI241\endcsname{\index{xiangrong tiaojian@相容条件（compatibility conditions）}}
\expandafter\def\csname liftmarkI242\endcsname{\index{xiangrong tiaojian@相容条件（compatibility conditions）}}
\expandafter\def\csname liftmarkI243\endcsname{\index{zhuangshi@装饰（decoration）}}
\expandafter\def\csname liftmarkI244\endcsname{\index{dai@代（generation）}}
\expandafter\def\csname liftmarkI245\endcsname{\index{ank daishu@$A_{n,K}$ 代数（$A_{n,K}$ algebra）}}
\expandafter\def\csname liftmarkI246\endcsname{\index{ank tongtai@$A_{n,K}$ 同态（$A_{n,K}$ homomorphism）}}
\expandafter\def\csname liftmarkI247\endcsname{\index{ank tonglun@$A_{n,K}$ 同伦（$A_{n,K}$ homotopy）}}
\expandafter\def\csname liftrangestartI248\endcsname{\index{lvguo ank daishu@滤过 $A_{n,K}$ 代数（filtered $A_{n,K}$ algebra）|(}}
\expandafter\def\csname liftrangeendI248\endcsname{\index{lvguo ank daishu@滤过 $A_{n,K}$ 代数（filtered $A_{n,K}$ algebra）|)}}
\expandafter\def\csname liftrangestartI249\endcsname{\index{lvguo ank daishu@滤过 $A_{n,K}$ 代数（filtered $A_{n,K}$ algebra）!dingyi@定义（definition）|(}}
\expandafter\def\csname liftrangeendI249\endcsname{\index{lvguo ank daishu@滤过 $A_{n,K}$ 代数（filtered $A_{n,K}$ algebra）!dingyi@定义（definition）|)}}
\expandafter\def\csname liftmarkI250\endcsname{\index{lvguo ank daishu@滤过 $A_{n,K}$ 代数（filtered $A_{n,K}$ algebra）!lvguo ank tongtai@滤过 $A_{n,K}$ 同态（filtered $A_{n,K}$ homomorphism）}}
\expandafter\def\csname liftmarkI251\endcsname{\index{lvguo ank daishu@滤过 $A_{n,K}$ 代数（filtered $A_{n,K}$ algebra）!lvguo ank tonglun@滤过 $A_{n,K}$ 同伦（filtered $A_{n,K}$ homotopy）}}
\expandafter\def\csname liftmarkI252\endcsname{\index{zhangai@障碍（obstruction）!kuozhang lvguo ank jiegou@扩张为滤过 $A_{n',K'}$ 结构（extending to a filtered $A_{n',K'}$ structure）}}
\expandafter\def\csname liftmarkI253\endcsname{\index{zhangai@障碍（obstruction）!kuozhang lvguo ank tongtai@扩张为滤过 $A_{n',K'}$ 同态（extending to a filtered $A_{n',K'}$ homomorphism）}}
\expandafter\def\csname liftmarkI254\endcsname{\index{ank daishu@$A_{n,K}$ 代数（$A_{n,K}$ algebra）!jihe shixian@几何实现（geometric realization）}}
\expandafter\def\csname liftmarkI255\endcsname{\index{lvguo ank daishu@滤过 $A_{n,K}$ 代数（filtered $A_{n,K}$ algebra）!lagelangri zi liuxing@拉格朗日子流形上的结构（of Lagrangian submanifolds）}}
\expandafter\def\csname liftmarkI256\endcsname{\index{ank tongtai@$A_{n,K}$ 同态（$A_{n,K}$ homomorphism）!jihe shixian@几何实现（geometric realization）}}
\expandafter\def\csname liftmarkI257\endcsname{\index{zhangai@障碍（obstruction）!kuozhang lvguo ank tonglun@扩张为滤过 $A_{n',K'}$ 同伦（to extending to a filtered $A_{n',K'}$ homotopy）}}
\expandafter\def\csname liftmarkI258\endcsname{\index{ank tonglun@$A_{n,K}$ 同伦（$A_{n,K}$ homotopy）!jihe shixian@几何实现（geometric realization）}}
\expandafter\def\csname liftmarkI259\endcsname{\index{shipei qiyi danxing@适配奇异单形（adapted singular simplex）}}
\expandafter\def\csname liftmarkI260\endcsname{\index{lengzhu poufen@棱柱剖分（prism decomposition）}}
\expandafter\def\csname liftmarkI261\endcsname{\index{incl@$\mathrm{Incl}$!jihe shixian@几何实现（geometric realization）}}
\expandafter\def\csname liftmarkI262\endcsname{\index{eval@$\mathrm{Eval}$}}
\expandafter\def\csname liftmarkI263\endcsname{\index{incl@$\mathrm{Incl}$}}
\expandafter\def\csname liftmarkI264\endcsname{\index{lingyi c moxing@$[0,1]\times C$ 的模型（model of $[0,1]\times C$）!awuqiong daishu@$A_\infty$ 代数（$A_\infty$ algebra，修订版）}}
\expandafter\def\csname liftmarkI265\endcsname{\index{kegouzao ji@可构造集（constructible set）}}
\expandafter\def\csname liftmarkI266\endcsname{\index{youxiang xitong@有向系统（directed system）!dengjia youxiang xitong@等价有向系统（of equivalences）}}
\expandafter\def\csname liftmarkI267\endcsname{\index{youxiang xitong@有向系统（directed system）!youxiang xitong taishi@有向系统的态射（morphism of）}}
\expandafter\def\csname liftmarkI268\endcsname{\index{tisheng@提升（promotion）}}
\expandafter\def\csname liftmarkI269\endcsname{\index{jiaohuan wucha tonglun@交换误差同伦（commuting error homotopy）!jiaohuan wucha tonglun taishi@态射的交换误差同伦（of morphisms）}}
\expandafter\def\csname liftmarkI270\endcsname{\index{erwei qujian moxing@$[0,1]^2\times C$ 的模型（model of $[0,1]^2\times C$）}}
\expandafter\def\csname liftmarkI271\endcsname{\index{jiaohuan wucha tonglun@交换误差同伦（commuting error homotopy）!jiaohuan wucha tonglun tisheng@提升的交换误差同伦（of the promotion）}}
\expandafter\def\csname liftrangestartI272\endcsname{\index{eval@$\mathrm{Eval}$|(}}
\expandafter\def\csname liftrangeendI272\endcsname{\index{eval@$\mathrm{Eval}$|)}}
\expandafter\def\csname liftrangestartI273\endcsname{\index{incl@$\mathrm{Incl}$|(}}
\expandafter\def\csname liftrangeendI273\endcsname{\index{incl@$\mathrm{Incl}$|)}}
\expandafter\def\csname liftmarkI274\endcsname{\index{jiaohuan wucha tonglun@交换误差同伦（commuting error homotopy）!jiaohuan wucha tonglun tonglun@同伦的交换误差同伦（of homotopies）}}
\expandafter\def\csname liftmarkI275\endcsname{\index{tonglun de tonglun@同伦的同伦（homotopy of homotopies）}}
\expandafter\def\csname liftmarkI276\endcsname{\index{incl@$\mathrm{Incl}$}}
\expandafter\def\csname liftmarkI277\endcsname{\index{tonglun de tonglun@同伦的同伦（homotopy of homotopies）}}
\expandafter\def\csname liftmarkI278\endcsname{\index{shipei qiyi danxing@适配奇异单形（adapted singular simplex）}}
\expandafter\def\csname liftmarkI279\endcsname{\index{lengzhu poufen@棱柱剖分（prism decomposition）}}
\expandafter\def\csname liftmarkI280\endcsname{\index{tonglun yaoyuan@同伦幺元（homotopy unit）}}
\expandafter\def\csname liftmarkI281\endcsname{\index{yiwang yingshe@遗忘映射（forgetful map）}}
\expandafter\def\csname liftmarkI282\endcsname{\index{yiwang yingshe@遗忘映射（forgetful map）}}
\expandafter\def\csname liftmarkI283\endcsname{\index{xiangrong tiaojian@相容条件（compatibility conditions）}}
\expandafter\def\csname liftmarkI284\endcsname{\index{xiangrong tiaojian@相容条件（compatibility conditions）}}
\expandafter\def\csname liftmarkI285\endcsname{\index{suanzi p@算子 $\mathfrak p$（operator $\mathfrak p$）}}
\expandafter\def\csname liftmarkI286\endcsname{\index{lwuqiong daishu@$L_\infty$ 代数（$L_\infty$ algebra）}}
\expandafter\def\csname liftmarkI287\endcsname{\index{lwuqiong tongtai@$L_\infty$ 同态（$L_\infty$ homomorphism）}}
\expandafter\def\csname liftmarkI288\endcsname{\index{fanxu@反序（opposite）}}
\expandafter\def\csname liftmarkI289\endcsname{\index{fanxu@反序（opposite）}}
\expandafter\def\csname liftmarkI290\endcsname{\index{daizhuang shu@带状树（ribbon tree）}}
\expandafter\def\csname liftmarkI291\endcsname{\index{daizhuang shu@带状树（ribbon tree）}}
\expandafter\def\csname liftmarkI292\endcsname{\index{suanzi r@算子 $\mathfrak r$（operator $\mathfrak r$）}}
\expandafter\def\csname liftmarkI293\endcsname{\index{awuqiong fanchou@$A_\infty$ 范畴（$A_\infty$ category）}}
\expandafter\def\csname liftmarkI294\endcsname{\index{awuqiong fanchou@$A_\infty$ 范畴（$A_\infty$ category）}}
\expandafter\def\csname liftmarkI295\endcsname{\index{lwuqiong daishu@$L_\infty$ 代数（$L_\infty$ algebra）}}
\expandafter\def\csname liftmarkI296\endcsname{\index{suanzi p@算子 $\mathfrak p$（operator $\mathfrak p$）!xunhuan bar fuxing@循环 bar 复形（cyclic bar complex）}}
\expandafter\def\csname liftmarkI297\endcsname{\index{dingxiang@定向（orientation）}}
\expandafter\def\csname liftmarkI298\endcsname{\index{dingxiang@定向（orientation）!wei quan chun yuanpan mokongjian@伪全纯圆盘模空间（moduli space of pseudo-holomorphic disc）}}
\expandafter\def\csname liftrangestartI299\endcsname{\index{dingxiang@定向（orientation）!bianjie@边界（boundary）|(}}
\expandafter\def\csname liftrangeendI299\endcsname{\index{dingxiang@定向（orientation）!bianjie@边界（boundary）|)}}
\expandafter\def\csname liftmarkI300\endcsname{\index{dingxiang@定向（orientation）!shang kongjian@商空间（quotient space）}}
\expandafter\def\csname liftmarkI301\endcsname{\index{dingxiang@定向（orientation）!xianwei ji@纤维积（fiber product）}}
\expandafter\def\csname liftmarkI302\endcsname{\index{dingxiang tonggou@定向同构（oriented isomorphism）}}
\expandafter\def\csname liftmarkI303\endcsname{\index{dingxiang@定向（orientation）!changzhi yingshe mokongjian@常值映射模空间（moduli space of constant maps）}}
\expandafter\def\csname liftmarkI304\endcsname{\index{dingxiang@定向（orientation）!lianjie guidao mokongjian@连接轨道模空间（moduli space of connecting orbits）}}
\expandafter\def\csname liftmarkI305\endcsname{\index{dingxiang@定向（orientation）!bote moersi qingxing@博特--莫尔斯情形（Bott--Morse cases）}}
\expandafter\def\csname liftmarkI306\endcsname{\index{dingxiang@定向（orientation）!dingxiang cong@定向丛（orientation bundle）}}
\expandafter\def\csname liftmarkI307\endcsname{\index{dingxiang@定向（orientation）!shijian youxu ji yu zhushi ji mokongjian@时间有序积与逐时积模空间（top and twp moduli spaces）}}
\expandafter\def\csname liftmarkI308\endcsname{\index{cangxi linyu@仓西邻域（Kuranishi neighborhood）}}
\expandafter\def\csname liftmarkI309\endcsname{\index{cangxi jiegou@仓西结构（Kuranishi structure）!lianghao zuobiao xi@良好坐标系（good coordinate system）}}
\expandafter\def\csname liftmarkI310\endcsname{\index{cangxi jiegou@仓西结构（Kuranishi structure）!zhangai cong@障碍丛（obstruction bundle）}}
\expandafter\def\csname liftmarkI311\endcsname{\index{cangxi jiegou@仓西结构（Kuranishi structure）!dingxiang@定向（orientation）}}
\expandafter\def\csname liftmarkI312\endcsname{\index{cangxi jiegou@仓西结构（Kuranishi structure）!qie cong@切丛（tangent bundle）}}
\expandafter\def\csname liftmarkI313\endcsname{\index{cangxi jiegou@仓西结构（Kuranishi structure）!xu jiben lian@虚基本链（virtual fundamental chain）}}
\expandafter\def\csname liftrangestartI314\endcsname{\index{duozhi jiemian@多值截面（multisection）|(}}
\expandafter\def\csname liftrangeendI314\endcsname{\index{duozhi jiemian@多值截面（multisection）|)}}
\expandafter\def\csname liftrangestartI315\endcsname{\index{duozhi jiemian@多值截面（multisection）!ketisheng duozhi jiemian@可提升多值截面（liftable multisection）|(}}
\expandafter\def\csname liftrangeendI315\endcsname{\index{duozhi jiemian@多值截面（multisection）!ketisheng duozhi jiemian@可提升多值截面（liftable multisection）|)}}
\expandafter\def\csname liftmarkI316\endcsname{\index{zhangai@障碍（obstruction）!zhangai cong@仓西结构的障碍丛（obstruction bundle of Kuranishi structure）}}
\expandafter\def\csname liftmarkI317\endcsname{\index{cangxi jiegou@仓西结构（Kuranishi structure）!xianwei ji@纤维积（fiber product）}}
\expandafter\def\csname liftmarkI318\endcsname{\index{cangxi jiegou@仓西结构（Kuranishi structure）!youxian qun zuoyong@有限群作用（finite group action）}}
\expandafter\def\csname liftmarkI319\endcsname{\index{awuqiong shuangmo tongtai@$A_\infty$ 双模同态（$A_\infty$ bimodule homomorphism）!c ruolvguo awuqiong shuangmo tongtai@$c$-弱滤过 $A_\infty$ 双模同态（$c$-weakly filtered $A_\infty$ bimodule homomorphism）}}
\expandafter\def\csname liftmarkI320\endcsname{\index{awuqiong daishu@$A_\infty$ 代数（$A_\infty$ algebra）!mo te wuzhangai@模 $T^E$ 无障碍（unobstructed modulo $T^E$）}}
\expandafter\def\csname liftmarkI321\endcsname{\index{awuqiong shuangmo@$A_\infty$ 双模（$A_\infty$ bimodule）!lahui@拉回（pull-back）}}
\expandafter\def\csname liftmarkI322\endcsname{\index{awuqiong shuangmo tongtai@$A_\infty$ 双模同态（$A_\infty$ bimodule homomorphism）!xingbian@形变（deformation）}}
\expandafter\def\csname liftmarkI323\endcsname{\index{awuqiong shuangmo tongtai@$A_\infty$ 双模同态（$A_\infty$ bimodule homomorphism）!mo t lambda lvguo awuqiong shuangmo tongtai@模 $T^\lambda$ 的滤过 $A_\infty$ 双模同态（filtered $A_\infty$ bimodule homomorphism modulo $T^\lambda$）}}
\expandafter\def\csname liftmarkI324\endcsname{\index{awuqiong shuangmo tongtai@$A_\infty$ 双模同态（$A_\infty$ bimodule homomorphism）!tonglun@同伦（homotopic）}}
\expandafter\def\csname liftmarkI325\endcsname{\index{awuqiong shuangmo tongtai@$A_\infty$ 双模同态（$A_\infty$ bimodule homomorphism）!lahui@拉回（pull-back）}}
\expandafter\def\csname liftmarkI326\endcsname{\index{awuqiong daishu tongtai@$A_\infty$（代数）同态（$A_\infty$ (algebra) homomorphism）!xingbian@形变（deformation）}}
\expandafter\def\csname liftmarkI327\endcsname{\index{awuqiong daishu tongtai@$A_\infty$（代数）同态（$A_\infty$ (algebra) homomorphism）!mo t lambda i lvguo awuqiong tongtai@模 $T^{\lambda_i}$ 的滤过 $A_\infty$ 同态（filtered $A_\infty$ homomorphisms modulo $T^{\lambda_i}$）}}
\expandafter\def\csname liftmarkI328\endcsname{\index{fencha lunzheng peibian lunzheng@分叉论证与配边论证（bifurcation versus cobordism）}}
\expandafter\def\csname liftmarkI329\endcsname{\index{bulk xingbian@bulk 形变（bulk deformations）!bulk xingbian fuluoer shang tongdiao@bulk 形变弗洛尔上同调（bulk-deformed Floer cohomology）}}
\expandafter\def\csname liftmarkI330\endcsname{\index{xingbian@形变（deformation）!awuqiong daishu tongtai@$A_\infty$ 代数同态的形变（of $A_\infty$ algebra homomorphisms）}}
\expandafter\def\csname liftmarkI331\endcsname{\index{xingbian@形变（deformation）!awuqiong shuangmo tongtai@$A_\infty$ 双模同态的形变（of $A_\infty$ bimodule homomorphisms）}}
\expandafter\def\csname liftmarkI332\endcsname{\index{tuihua danxing@退化单形（degenerate simplex）}}
\expandafter\def\csname liftmarkI333\endcsname{\index{delamu@德拉姆（de Rham）!xu jiben liande delamu biaoshi@虚基本链的德拉姆表示（presentation of virtual fundamental chain）}}
\expandafter\def\csname liftmarkI334\endcsname{\index{d g c f z@微分分次 c.f.z.（d.g.c.f.z.）}}
\expandafter\def\csname liftmarkI335\endcsname{\index{lvguo lwuqiong shuangmo@滤过 $L_\infty$ 双模（filtered $L_\infty$ bimodule）}}
\expandafter\def\csname liftmarkI336\endcsname{\index{guifan dengjia@规范等价（gauge equivalence）!g guifan dengjia@$G$-规范等价（$G$-gauge equivalent）}}
\expandafter\def\csname liftmarkI337\endcsname{\index{g ank jiegou@$G$-$A_{n;K_1,K_0}$ 结构（$G$-$A_{n;K_1,K_0}$-structure）}}
\expandafter\def\csname liftmarkI338\endcsname{\index{hochschild shangtongdiao@Hochschild 上同调（Hochschild cohomology）!awuqiong shuangmo@$A_\infty$ 双模的（of $A_\infty$ bimodule）}}
\expandafter\def\csname liftmarkI339\endcsname{\index{xiangdui bianjie tonglun@相对于边界的同伦（homotopy relative to the boundary）}}
\expandafter\def\csname liftmarkI340\endcsname{\index{diedai gongshi@迭代公式（iteration formula）}}
\expandafter\def\csname liftmarkI341\endcsname{\index{lwuqiong shuangmo@$L_\infty$ 双模（$L_\infty$ bimodule）}}
\expandafter\def\csname liftmarkI342\endcsname{\index{jixiao moxing@极小模型（minimal model）}}
\expandafter\def\csname liftmarkI343\endcsname{\index{zuixiao lujing@最小路径（minimal path）}}
\expandafter\def\csname liftrangestartI344\endcsname{\index{lingyi c moxing@$[0,1]\times C$ 的模型（model of $[0,1]\times C$）!awuqiong shuangmo@$A_\infty$ 双模（$A_\infty$ bimodule）|(}}
\expandafter\def\csname liftrangeendI344\endcsname{\index{lingyi c moxing@$[0,1]\times C$ 的模型（model of $[0,1]\times C$）!awuqiong shuangmo@$A_\infty$ 双模（$A_\infty$ bimodule）|)}}
\expandafter\def\csname liftmarkI345\endcsname{\index{youbian wending yingshe mo kongjian@有边稳定映射模空间（moduli space of bordered stable maps）!xu jiben lian@虚基本链（virtual fundamental chain）}}
\expandafter\def\csname liftmarkI346\endcsname{\index{nuoweikefu tongdiao@诺维科夫同调（Novikov homology）}}
\expandafter\def\csname liftmarkI347\endcsname{\index{nuoweikefu huan@诺维科夫环（Novikov ring）!hamidun tonghen bubianxing@哈密顿同痕下的无关性（independence under Hamiltonian isotopy）}}
\expandafter\def\csname liftrangestartI348\endcsname{\index{zhangai@障碍（obstruction）!fuluoer shang tongdiao zhangai lei@弗洛尔上同调的障碍类（obstruction classes for Floer cohomology）|(}}
\expandafter\def\csname liftrangeendI348\endcsname{\index{zhangai@障碍（obstruction）!fuluoer shang tongdiao zhangai lei@弗洛尔上同调的障碍类（obstruction classes for Floer cohomology）|)}}
\expandafter\def\csname liftmarkI349\endcsname{\index{zhangai@障碍（obstruction）!fuluoer shang tongdiao zhangai lei@弗洛尔上同调的障碍类（obstruction classes for Floer cohomology）}}
\expandafter\def\csname liftmarkI350\endcsname{\index{zhangai@障碍（obstruction）!gouzao lvguo awuqiong shuangmo tongtai@构造滤过 $A_\infty$ 双模同态的障碍（to constructing filtered $A_\infty$ bimodule homomorphisms）}}
\expandafter\def\csname liftmarkI351\endcsname{\index{zhangai@障碍（obstruction）!kuozhang ak jia yi tongtai@扩张为 $A_{K+1}$ 同态的障碍（to extending to $A_{K+1}$ homomorphism）}}
\expandafter\def\csname liftmarkI352\endcsname{\index{zhangai@障碍（obstruction）!kuozhang mo t lambda i jia yi@扩张到模 $T^{\lambda_{i+1}}$ 的障碍（to extending to modulo $T^{\lambda_{i+1}}$）}}
\expandafter\def\csname liftmarkI353\endcsname{\index{zhangai@障碍（obstruction）!ruowuzhangai@弱无障碍（weakly unobstructed）}}
\expandafter\def\csname liftmarkI354\endcsname{\index{zhangai@障碍（obstruction）!bulk xingbian hou ruowuzhangai@经 bulk 形变后（弱）无障碍（(weakly) unobstructed after bulk deformations）}}
\expandafter\def\csname liftmarkI355\endcsname{\index{suanzi p@算子 $\mathfrak p$（operator $\mathfrak p$）!tuiguang@推广（generalization）}}
\expandafter\def\csname liftmarkI356\endcsname{\index{suanzi q@算子 $\mathfrak q$（operator $\mathfrak q$）!gouzao gailue@构造概述（outline of construction）}}
\expandafter\def\csname liftrangestartI357\endcsname{\index{suanzi q@算子 $\mathfrak q$（operator $\mathfrak q$）!dingyi@定义（definition）|(}}
\expandafter\def\csname liftrangeendI357\endcsname{\index{suanzi q@算子 $\mathfrak q$（operator $\mathfrak q$）!dingyi@定义（definition）|)}}
\expandafter\def\csname liftrangestartI358\endcsname{\index{suanzi q@算子 $\mathfrak q$（operator $\mathfrak q$）!dingyi@定义（definition）|(}}
\expandafter\def\csname liftrangeendI358\endcsname{\index{suanzi q@算子 $\mathfrak q$（operator $\mathfrak q$）!dingyi@定义（definition）|)}}
\expandafter\def\csname liftmarkI359\endcsname{\index{suanzi q@算子 $\mathfrak q$（operator $\mathfrak q$）!dingyi@定义（definition）}}
\expandafter\def\csname liftmarkI360\endcsname{\index{dingxiang@定向（orientation）!lianxu raodong zu@连续扰动族（continuous family of perturbations）}}
\expandafter\def\csname liftmarkI361\endcsname{\index{dingxiang@定向（orientation）!tonglun yaoyuan@同伦幺元（homotopy unit）}}
\expandafter\def\csname liftmarkI362\endcsname{\index{poly@$\operatorname{Poly}$}}
\expandafter\def\csname liftmarkI363\endcsname{\index{youli lagelangri zi liuxing@有理拉格朗日子流形（rational Lagrangian submanifold）}}
\expandafter\def\csname liftrangestartI364\endcsname{\index{youli lagelangri zi liuxing@有理拉格朗日子流形（rational Lagrangian submanifold）|(}}
\expandafter\def\csname liftrangeendI364\endcsname{\index{youli lagelangri zi liuxing@有理拉格朗日子流形（rational Lagrangian submanifold）|)}}
\expandafter\def\csname liftmarkI365\endcsname{\index{xiangdui spin@相对 spin（relatively spin）!jiegou@结构（structure）}}
\expandafter\def\csname liftmarkI366\endcsname{\index{daijubu xishu qiyi lian@带局部系数的奇异链（singular chains with local coefficients）}}
\expandafter\def\csname liftmarkI367\endcsname{\index{pu xulie@谱序列（spectral sequence）}}
\expandafter\def\csname liftmarkI368\endcsname{\index{pu xulie@谱序列（spectral sequence）!guanlian lagelangri zi liuxing@与拉格朗日子流形关联（associated Lagrangian submanifolds）}}
\expandafter\def\csname liftmarkI369\endcsname{\index{pu xulie@谱序列（spectral sequence）!shoulianxing@收敛性（convergence）}}
\expandafter\def\csname liftmarkI370\endcsname{\index{pu xulie@谱序列（spectral sequence）!wu longgen pu xulie@吴龙根谱序列（Oh's spectral sequence）}}
\expandafter\def\csname liftmarkI371\endcsname{\index{yange shoulian mijishu huan@严格收敛幂级数环（strictly convergent power series ring）}}
\expandafter\def\csname liftmarkI372\endcsname{\index{qiang lianxu@强连续（strongly continuous）}}
\expandafter\def\csname liftmarkI373\endcsname{\index{chao jiaohuanzi@超交换子（super commutator）}}
\expandafter\def\csname liftmarkI374\endcsname{\index{youbian wending yingshe mo kongjian tapu@有边稳定映射模空间的拓扑（topology of the moduli space of bordered stable maps）}}
\expandafter\def\csname liftmarkI375\endcsname{\index{er fanchou@2-范畴（2-category）}}
\expandafter\def\csname liftrangestartI376\endcsname{\index{mo te wuzhangai@模 $T^E$ 无障碍（unobstructed modulo $T^E$）|(}}
\expandafter\def\csname liftrangeendI376\endcsname{\index{mo te wuzhangai@模 $T^E$ 无障碍（unobstructed modulo $T^E$）|)}}
\expandafter\def\csname liftmarkI377\endcsname{\index{xu jiben shanglian@虚基本上链（virtual fundamental cochain）}}
\expandafter\def\csname liftmarkI378\endcsname{\index{ruolvguo awuqiong shuangmo tongtai@弱滤过 $A_\infty$ 双模同态（weakly filtered $A_\infty$ bimodule homomorphism）}}
\expandafter\def\csname liftmarkI379\endcsname{\index{ruolvguo g ank shuangmo tongtai@弱滤过 $G'$-$A_{n;K_1,K_0}$ 双模同态（weakly filtered $G'$-$A_{n;K_1,K_0}$ bimodule homomorphism）}}
\expandafter\def\csname liftmarkI380\endcsname{\index{ruoyouxian@弱有限（weakly finite）}}
\expandafter\def\csname liftmarkI381\endcsname{\index{ruojinmo@弱浸没（weakly submersive）}}
\expandafter\def\csname liftrangestartI382\endcsname{\index{ruojinmo@弱浸没（weakly submersive）|(}}
\expandafter\def\csname liftrangeendI382\endcsname{\index{ruojinmo@弱浸没（weakly submersive）|)}}
\expandafter\def\csname liftmarkI383\endcsname{\index{whitehead dingli@Whitehead 定理（Whitehead theorem）!lwuqiong daishu@对 $L_\infty$ 代数（for $L_\infty$ algebras）}}
\expandafter\def\csname liftrangestartI384\endcsname{\index{bote moersi lilun@博特--莫尔斯理论（Bott--Morse theory）|(}}
\expandafter\def\csname liftrangeendI384\endcsname{\index{bote moersi lilun@博特--莫尔斯理论（Bott--Morse theory）|)}}
\expandafter\def\csname liftrangestartI385\endcsname{\index{xunhuan@循环（Cyclic）!xunhuan duichengxing@循环对称性（symmetry）|(}}
\expandafter\def\csname liftrangeendI385\endcsname{\index{xunhuan@循环（Cyclic）!xunhuan duichengxing@循环对称性（symmetry）|)}}
\expandafter\def\csname liftrangestartI386\endcsname{\index{incl@$\mathrm{Incl}$!jihe shixian@几何实现（geometric realization）|(}}
\expandafter\def\csname liftrangeendI386\endcsname{\index{incl@$\mathrm{Incl}$!jihe shixian@几何实现（geometric realization）|)}}
\expandafter\def\csname liftmarkI387\endcsname{\index{suanzi q@算子 $\mathfrak q$（operator $\mathfrak q$）}}
\expandafter\def\csname liftmarkI388\endcsname{\index{suanzi p@算子 $\mathfrak p$（operator $\mathfrak p$）}}
\expandafter\def\csname liftmarkI389\endcsname{\index{yiwang yingshe@遗忘映射（forgetful map）}}

\DeclareRobustCommand{\liftIndex}[1]{\csname liftmark#1\endcsname}
\DeclareRobustCommand{\liftIndexRangeStart}[1]{\csname liftrangestart#1\endcsname}
\DeclareRobustCommand{\liftIndexRangeEnd}[1]{\csname liftrangeend#1\endcsname}


\makeatletter
\renewenvironment{theindex}{%
  \chapter*{\indexname}%
  \addcontentsline{toc}{chapter}{\indexname}%
  \markboth{\indexname}{\indexname}%
  \begingroup
  \footnotesize
  \setlength{\parindent}{0pt}%
  \setlength{\parskip}{0pt}%
  \rightskip=0pt plus 1fil\relax
  \setlength{\columnsep}{1.5em}%
  \noindent\textbf{凡例}\quad 本索引按中文无声调拼音排序；括号中保留英文源词，页码使用本译本正文印刷页码。人名对照另见“双向人名检索表”。%
  \par\vspace{0.55\baselineskip}%
  \begin{multicols}{2}%
  \raggedcolumns
  \let\item\@idxitem
}{%
  \end{multicols}%
  \endgroup
}
\renewcommand{\@idxitem}{\par\hangindent=1em\hangafter=1\relax}
\renewcommand{\subitem}{\par\hangindent=2.8em\hangafter=1\hspace*{1.4em}}
\makeatother

\pagestyle{fancy}
\fancyhf{}
\fancyhead[LE]{\small\thepage\quad\nouppercase{\leftmark}}
\fancyhead[RO]{\small\nouppercase{\rightmark}\quad\thepage}
\renewcommand{\headrulewidth}{0.4pt}

\makeatletter
\newcommand{\liftstatementhead}[2]{%
  \par\addvspace{\topsep}%
  \noindent\phantomsection%
  \protected@edef\@currentlabel{#2}%
  \textbf{#1~#2.}\enspace\ignorespaces
}
\makeatother
\newcommand{\liftstatementend}{\par\addvspace{\topsep}}

\newenvironment{statement}[2]
  {\liftstatementhead{#1}{#2}}
  {\liftstatementend}
\newenvironment{theorem}[1]
  {\begin{statement}{定理}{#1}\itshape}
  {\end{statement}}
\newenvironment{proposition}[1]
  {\begin{statement}{命题}{#1}\itshape}
  {\end{statement}}
\newenvironment{lemma}[1]
  {\begin{statement}{引理}{#1}\itshape}
  {\end{statement}}
\newenvironment{corollary}[1]
  {\begin{statement}{推论}{#1}\itshape}
  {\end{statement}}
\newenvironment{conjecture}[1]
  {\begin{statement}{猜想}{#1}\itshape}
  {\end{statement}}
\newenvironment{definition}[1]
  {\begin{statement}{定义}{#1}}
  {\end{statement}}
\newenvironment{example}[1]
  {\begin{statement}{例}{#1}}
  {\end{statement}}
\newenvironment{exercise}[1]
  {\begin{statement}{习题}{#1}}
  {\end{statement}}
\newenvironment{remark}[1]
  {\begin{statement}{注}{#1}}
  {\end{statement}}


\usepackage{booktabs,longtable,multicol}
\newcommand{\BookGlossEntry}[2]{%
  \par\noindent\hangindent=1em\hangafter=1 #1\quad #2\par}
\newcommand{\BookSourceIndexEntry}[2]{%
  \par\begingroup\leftskip=#1em\relax\hangindent=1em\relax
  \noindent #2\par\endgroup}
\makeatletter
\newif\ifBookIndexing
\BookIndexingfalse
\newcommand{\BookBoldPage}[1]{\textbf{\hyperpage{#1}}}
\DeclareRobustCommand{\BookRecordIndex}[1]{%
  \ifBookIndexing
    \expandafter\index\expandafter{\csname BookIndex@#1\endcsname|hyperpage}%
  \fi}
\DeclareRobustCommand{\BookIndexDefinition}[1]{%
  \ifBookIndexing
    \expandafter\index\expandafter{\csname BookIndex@#1\endcsname|BookBoldPage}%
  \fi}
\DeclareRobustCommand{\BookTerm}[2]{%
  \BookRecordIndex{#1}%
  \begingroup\let\BookRecordIndex\@gobble #2\endgroup}
\AddToHook{cmd/@outputpage/before}{\let\BookRecordIndex\@gobble}
\pdfstringdefDisableCommands{%
  \def\BookTerm#1#2{#2}%
  \let\BookRecordIndex\@gobble
  \let\BookIndexDefinition\@gobble}
\makeatother

\let\BookOriginalIndex\index
\renewcommand{\index}[1]{\ifBookIndexing\BookOriginalIndex{#1}\fi}
